% Subsection: Closure Under Finite Unions and Intersections  (notes stub; content authored later)
\subsection{Closure Under Finite Unions and Intersections}

\begin{remark*}[Exposition]
Finite union and finite intersection closure are the basic finite stability
conditions for a set system.  They say that a family can tolerate finitely
many join-like or meet-like combinations without leaving the family.
\end{remark*}

\begin{remark*}[Exposition]
Many finite closure requirements are equivalent to their two-input versions
provided the family is nonempty in the relevant way.  For example, repeated
use of binary union gives finite union closure.  The finite formulation is
nevertheless the more stable statement because it records the full operation
the family is expected to support.
\end{remark*}

\begin{definition}[Closed Under Finite Unions]\label{def:closed-under-finite-unions}
Let \(X\) be a set and let \(\mathcal{F}\subseteq\mathcal{P}(X)\).  The family
\(\mathcal{F}\) is \emph{closed under finite unions} if for every finite list
\(A_1,\dots,A_n\in\mathcal{F}\), with \(n\geq 1\), one has
\[
  \bigcup_{k=1}^{n} A_k \in \mathcal{F}.
\]
\end{definition}

\begin{remark*}[Standard quantified statement]
\[
  \forall n\in\mathbb{N}_{\geq 1}\,
  \forall A_1,\dots,A_n\in\mathcal{F},\qquad
  \bigcup_{k=1}^{n}A_k\in\mathcal{F}.
\]
\end{remark*}

\begin{remark*}[Predicate reading]
\[
\begin{aligned}
&C_{\mathrm{fin}\cup}
=\mathsf{ClosureContext}(X,\mathcal{F}^{<\omega}_{\geq 1},\mathcal{F},\bigcup_{\mathrm{fin}}),\\
&\operatorname{ClosedUnder}(\bigcup_{\mathrm{fin}},\mathcal{F},C_{\mathrm{fin}\cup}).
\end{aligned}
\]
\end{remark*}

\begin{remark*}[Interpretation]
Closure under finite unions says that combining finitely many admitted sets by
``or'' gives another admitted set.
\end{remark*}

\begin{dependencies}
\begin{itemize}
  \item \hyperref[def:closed-under-operation]{Closed Under an Operation}
  \item \hyperref[def:finite-operation-on-family]{Finite Operation on a Family}
  \item \hyperref[def:union]{Union}
  \item \hyperref[def:finite-infinite]{Finite and Infinite Sets}
\end{itemize}
\end{dependencies}

\begin{definition}[Closed Under Finite Intersections]\label{def:closed-under-finite-intersections}
Let \(X\) be a set and let \(\mathcal{F}\subseteq\mathcal{P}(X)\).  The family
\(\mathcal{F}\) is \emph{closed under finite intersections} if for every
finite list \(A_1,\dots,A_n\in\mathcal{F}\), with \(n\geq 1\), one has
\[
  \bigcap_{k=1}^{n} A_k \in \mathcal{F}.
\]
\end{definition}

\begin{remark*}[Standard quantified statement]
\[
  \forall n\in\mathbb{N}_{\geq 1}\,
  \forall A_1,\dots,A_n\in\mathcal{F},\qquad
  \bigcap_{k=1}^{n}A_k\in\mathcal{F}.
\]
\end{remark*}

\begin{remark*}[Predicate reading]
\[
\begin{aligned}
&C_{\mathrm{fin}\cap}
=\mathsf{ClosureContext}(X,\mathcal{F}^{<\omega}_{\geq 1},\mathcal{F},\bigcap_{\mathrm{fin}}),\\
&\operatorname{ClosedUnder}(\bigcap_{\mathrm{fin}},\mathcal{F},C_{\mathrm{fin}\cap}).
\end{aligned}
\]
\end{remark*}

\begin{remark*}[Interpretation]
Closure under finite intersections says that combining finitely many admitted
sets by ``and'' gives another admitted set.
\end{remark*}

\begin{dependencies}
\begin{itemize}
  \item \hyperref[def:closed-under-operation]{Closed Under an Operation}
  \item \hyperref[def:finite-operation-on-family]{Finite Operation on a Family}
  \item \hyperref[def:intersection]{Intersection}
  \item \hyperref[def:finite-infinite]{Finite and Infinite Sets}
\end{itemize}
\end{dependencies}
